%% Chapter 1 — The Unverifiability Problem: Historical and Mathematical

\chapter{The Unverifiability Problem: Historical and Mathematical}
\label{chap:unverifiability}

\epigraph{%
  ``A theory which is not refutable by any conceivable event is non-scientific.''
}{--- Karl Popper, \textit{Conjectures and Refutations}, 1963}

\section{The Epistemological Gap in Classical AI}
\label{sec:epistemological-gap}

When a civil engineer stamps a load calculation, the profession does not trust the
engineer's assurances—it trusts the \emph{methodology}: a documented procedure
whose inputs and outputs can be independently recomputed, and whose failure modes are
known in advance.  When a pharmacologist submits a compound for regulatory review,
the submission package includes not only efficacy data but a complete specification
of the assay conditions under which those data were produced, so that any competent
laboratory can reproduce or refute the claim.  In both domains, \emph{operational
verifiability}—the existence of a concise certificate that a third party can check
in polynomial time—is the precondition for professional legitimacy.

Classical AI systems of the 1970s and 1980s possessed no such certificate.
MYCIN \cite{shortliffe1976} demonstrated that a rule-based consultation system
could match or exceed the diagnostic accuracy of specialist physicians on
bacteremia cases, yet the 1976 binary encoding of MYCIN was not operationally
verifiable in any professional sense.  A board evaluating MYCIN faced an
epistemological asymmetry: the system could be queried on test cases, and its
outputs could be compared with expert consensus, but no observer could distinguish
MYCIN's causal rule-firing machinery from a finite lookup table that happened to
agree with MYCIN on every tested input.

This asymmetry is not merely philosophical.  Consider a lookup table $L$ indexed
by the symptom vectors appearing in the published MYCIN evaluation suite.  For
every case in that suite, $L$ can reproduce MYCIN's recommended antibiotic and
confidence factor exactly.  Yet $L$ encodes zero causal knowledge; it will
produce arbitrary outputs on any input outside its index.  Because the 1976
evaluation did not exhaustively sample the input space—indeed, could not, given
that the space of symptom combinations is combinatorially vast—no test run on that
evaluation could distinguish MYCIN from $L$.  The binary was, in the precise sense
defined below, \emph{operationally non-falsifiable}.

\begin{definition}[Operational Verifiability]
\label{def:op-verifiability}
A function $f : X \to Y$ is \emph{operationally verifiable} with respect to a
specification $S \subseteq X \times Y$ if there exists a polynomial-time computable
predicate $\pi(x, y)$ such that
\[
  (x, y) \in S \iff \pi(x, f(x)) = 1
\]
for all $x \in X$.  The predicate $\pi$ is called a \emph{verification certificate}
for $f$ against $S$.
\end{definition}

In practice, $\pi$ is realized by the combination of (a) a cryptographic receipt
binding the input to the output and (b) a conformance check that the output
satisfies the structural invariants declared in $S$.

\begin{definition}[Operational Falsifiability]
\label{def:op-falsifiability}
A function $f : X \to Y$ is \emph{operationally falsifiable} with respect to a
specification $S$ if there exists a finite test suite $T \subset X$ such that for
every $g : X \to Y$ with $g \neq f$ (as functions), there exists some $x \in T$
with $(x, g(x)) \notin S$ or $g(x) \neq f(x)$.  Equivalently, $T$
\emph{separates} $f$ from every other implementation in the implementation space
$Y^X$.
\end{definition}

\begin{remark}
Popper's falsifiability criterion \cite{popper1963} applies to \emph{theories}:
a scientific theory must forbid at least one conceivable observation.  Applied to a
\emph{running system}, the criterion sharpens: the system must, via its
receipts and invariant checks, forbid at least one conceivable output for every
input.  A system that accepts any output—because it emits no receipt and enforces
no invariant—is not a scientific instrument; it is a black box.  The distinction
between a theory and a running system is precisely that the running system produces
concrete inputs and outputs on which falsifiability can be directly tested, not
merely conjectured.
\end{remark}

The MYCIN non-falsifiability argument can now be stated precisely.  Let $X_{\text{MYCIN}}$
be the space of patient symptom vectors and $Y_{\text{MYCIN}}$ be the space of
antibiotic recommendations with associated confidence factors.  The 1976 binary
produces outputs via the Certainty Factor calculus---the rule
$\CF(H,E) = \CF_{\mathrm{old}} + \CF_{\mathrm{new}}(1-\CF_{\mathrm{old}})$
for evidence with the same sign, as defined by \citet{shortliffe1976}---but the binary itself emits no certificate
binding its inputs to its outputs, and no invariant check that would distinguish
rule-firing from table lookup.  The evaluation suite $T_{\text{eval}}$ used in the
published study was finite and small relative to $\lvert X_{\text{MYCIN}} \rvert$.
Therefore, by Definition~\ref{def:op-falsifiability}, MYCIN was not operationally
falsifiable against $T_{\text{eval}}$: the lookup table $L$ described above passes
every test in $T_{\text{eval}}$ while violating the declared causal specification.

This is not a criticism of Shortliffe's achievement—MYCIN was a landmark.  It is
a diagnosis of a structural gap that persisted in AI system evaluation for decades
and that the Periodic Table of Reason is designed to close.

\section{The wasm4pm Validation Framework}
\label{sec:wasm4pm-framework}

The \textsc{wasm4pm} platform operationalizes the requirements implicit in
Definitions~\ref{def:op-verifiability} and~\ref{def:op-falsifiability} through
five independently necessary components.

\begin{enumerate}[label=(\roman*)]

\item \textbf{Deterministic WASM boundary.}  Each breed is compiled to a
WebAssembly module with a fixed binary interface.  WASM's memory isolation and
deterministic execution semantics guarantee that, given the same input, the module
produces bit-identical output across all conforming WASM runtimes.  Alone, this
component establishes reproducibility but not verifiability: identical outputs on
identical inputs do not prove that the computation follows the declared algorithm.

\item \textbf{Structural invariant suite.}  The specification $S_i$ for breed $i$
is encoded as a set of executable predicates (the ``proof gate'' suite) that are
evaluated against every output at the WASM boundary.  Alone, structural invariants
establish that outputs are \emph{well-formed} but do not bind them to the specific
inputs that caused them: a cached or replayed output would pass invariant checks.

\item \textbf{BLAKE3 receipt chain.}  The function $\blake$ computes a
cryptographic digest of the input and output pair, forming a tamper-evident receipt
that is appended to the receipt chain in \texttt{.wasm4pm/receipts/latest.json}.
Alone, the receipt chain establishes causal binding between a specific input and a
specific output but does not guarantee that the output is correct: a receipt can be
computed for any $(x, y)$ pair, including one where $y$ is wrong.

\item \textbf{Process-mining conformance check.}  The OTEL spans emitted during
execution are converted to an OCEL event log and replayed against the declared
process model using native Rust DFA replay via \texttt{validate\_ocel\_alignment()} (\texttt{ocel.rs}) against the declared lifecycle OCPN---deterministic, reproducible, no external process-mining framework.  This check detects deviations between the
declared algorithm sequence and the actual runtime sequence—catching, for example,
a breed that skips a mandatory inference step.  Alone, conformance checking
establishes that the right steps occurred in the right order but says nothing about
whether the outputs of those steps are numerically correct.

\item \textbf{Input hash at the WASM boundary.}  Beginning with the current
implementation, \texttt{cognition\_run()} in
\texttt{crates/wasm4pm-cognition/src/wasm.rs} computes
$\text{input\_hash} = \blake(\text{input\_json})$ and includes it in the
\texttt{ContractResult} receipt.  This component closes the gap left by (iii)
alone: because the input hash is computed \emph{inside} the WASM module before
any output is produced, an adversary cannot substitute a precomputed output without
invalidating the input hash.  Alone, however, this component does not establish
conformance of the algorithm sequence.
\end{enumerate}

The necessity of all five components follows from the one-sentence arguments above:
each component closes exactly one attack surface while leaving the others open.
Only their conjunction achieves operational falsifiability in the sense of
Definition~\ref{def:op-falsifiability}: a test suite $T$ that distinguishes a
conforming breed from any non-conforming alternative must be able to (a) reproduce
the computation, (b) check output structure, (c) verify causal binding, (d) verify
algorithmic sequence, and (e) verify input integrity.

\begin{remark}
The five-component framework is the \emph{minimal} conjunction.  One could add
further components—for example, formal verification of the WASM binary against a
mathematical specification—but the five listed are necessary in the sense that
dropping any one of them leaves a constructible counterexample (a non-conforming
implementation that passes all remaining checks).
\end{remark}

\section{Summary}
\label{sec:ch1-summary}

This chapter established the epistemological gap between classical AI evaluation
and the professional certification standards of engineering and pharmacology.  The
gap was made precise via Definitions~\ref{def:op-verifiability}
and~\ref{def:op-falsifiability}, and illustrated concretely by the MYCIN
non-falsifiability argument.  The five-component \textsc{wasm4pm} validation
framework was presented as the minimal conjunction that closes this gap.  Chapter~2
formalizes the categorical structure through which these components are organized:
the \emph{breed category}, in which each breed is an object with a declared
signature and each conformance-preserving transformation is a morphism.
